% Lista de siglas. Cada sigla con su significado; los términos en inglés en
% cursiva. Mantener en orden alfabético; misma lista en PI - Glosario.md.
\chapter*{Lista de siglas}
\addcontentsline{toc}{chapter}{Lista de siglas}
\markboth{Lista de siglas}{Lista de siglas}

{\small
\begin{longtable}{|>{\raggedright\arraybackslash}p{0.14\textwidth}%
  |>{\raggedright\arraybackslash}p{\dimexpr0.86\textwidth-4\tabcolsep-3\arrayrulewidth\relax}|}
\hline \encabezado \textbf{Sigla} & \textbf{Significado} \\ \hline
\endfirsthead
\hline \encabezado \textbf{Sigla} & \textbf{Significado} \\ \hline
\endhead
ADC & Conversor analógico-digital (\emph{analog-to-digital converter}). \\ \hline
ADU & Unidad del conversor analógico-digital (\emph{analog-to-digital unit}). \\ \hline
AMP & Multiprocesamiento asimétrico (\emph{Asymmetric Multi-processing}). \\ \hline
ANSI & \emph{American National Standards Institute}. \\ \hline
API & Interfaz de programación de aplicaciones (\emph{Application Programming Interface}). \\ \hline
APP & Aplicación cliente del sistema. \\ \hline
APU & Unidad de procesamiento de aplicaciones del Zynq (\emph{Application Processing Unit}). \\ \hline
ARN & Autoridad Regulatoria Nuclear. \\ \hline
ASG & Generador de señales arbitrarias (\emph{Arbitrary Signal Generator}). \\ \hline
ASIC & Circuito integrado de aplicación específica (\emph{Application-Specific Integrated Circuit}). \\ \hline
AXI & \emph{Advanced eXtensible Interface}. \\ \hline
BRAM & Bloque de memoria RAM de la FPGA (\emph{Block RAM}). \\ \hline
CDC & Cruce de dominios de reloj (\emph{clock domain crossing}). \\ \hline
CE & Habilitación de reloj (\emph{clock enable}). \\ \hline
CLB & Bloque lógico configurable (\emph{Configurable Logic Block}). \\ \hline
CMOS & Semiconductor complementario de óxido metálico (\emph{Complementary Metal-Oxide-Semiconductor}). \\ \hline
CMT & Bloque de gestión de reloj (\emph{Clock Management Tile}). \\ \hline
CPU & Unidad central de procesamiento (\emph{Central Processing Unit}). \\ \hline
CR-RC & Conformador pasaaltos (CR) seguido de pasabajos (RC). \\ \hline
DAC & Conversor digital-analógico (\emph{digital-to-analog converter}). \\ \hline
DC & Corriente continua (\emph{direct current}). \\ \hline
DDA & Distribución diferencial de alturas. \\ \hline
DDR & Memoria de doble tasa de datos (\emph{Double Data Rate}). \\ \hline
DMA & Acceso directo a memoria (\emph{Direct Memory Access}). \\ \hline
DNL & No linealidad diferencial (\emph{differential non-linearity}). \\ \hline
DSP & Procesamiento digital de señales (\emph{digital signal processing}); también el bloque DSP de la FPGA. \\ \hline
EEPROM & Memoria de sólo lectura programable y borrable eléctricamente. \\ \hline
EOF & Fin de trama (\emph{end of frame}). \\ \hline
FIFO & Cola en la que el primero en entrar es el primero en salir (\emph{First-In First-Out}). \\ \hline
FIR & Respuesta finita al impulso (\emph{Finite Impulse Response}). \\ \hline
FPGA & Arreglo de compuertas programable en campo (\emph{Field-Programmable Gate Array}). \\ \hline
FS & Fondo de escala (\emph{full scale}). \\ \hline
FSM & Máquina de estados finitos (\emph{Finite State Machine}). \\ \hline
FWHM & Ancho a media altura (\emph{full width at half maximum}). \\ \hline
GIL & Cerrojo global del intérprete de Python (\emph{Global Interpreter Lock}). \\ \hline
GP & Puerto AXI de propósito general del Zynq (\emph{General Purpose}). \\ \hline
GPIO & Entradas y salidas de propósito general (\emph{General Purpose Input/Output}). \\ \hline
GUI & Interfaz gráfica de usuario (\emph{Graphical User Interface}). \\ \hline
HDF5 & \emph{Hierarchical Data Format}, versión 5. \\ \hline
HP & Puerto AXI de alto rendimiento del Zynq (\emph{High Performance}). \\ \hline
HV, LV & Rango de entrada de alta y de baja tensión (\emph{high voltage}, \emph{low voltage}). \\ \hline
HW & Hardware. \\ \hline
I2C & \emph{Inter-Integrated Circuit}. \\ \hline
IAEA & Organismo Internacional de Energía Atómica (\emph{International Atomic Energy Agency}). \\ \hline
IEC & Comisión Electrotécnica Internacional (\emph{International Electrotechnical Commission}). \\ \hline
IEEE & \emph{Institute of Electrical and Electronics Engineers}. \\ \hline
IIR & Respuesta infinita al impulso (\emph{Infinite Impulse Response}). \\ \hline
INL & No linealidad integral (\emph{integral non-linearity}). \\ \hline
IP & Bloque de propiedad intelectual (\emph{Intellectual Property}). \\ \hline
ISA & Arquitectura del conjunto de instrucciones (\emph{Instruction Set Architecture}). \\ \hline
JSON & \emph{JavaScript Object Notation}. \\ \hline
LH5 & Formato de archivo de la colaboración LEGEND, basado en HDF5. \\ \hline
LLD & Discriminador de nivel inferior (\emph{lower level discriminator}). \\ \hline
LSB & Bit menos significativo (\emph{least significant bit}). \\ \hline
LUT & Tabla de consulta (\emph{Look-Up Table}). \\ \hline
MCA & Analizador multicanal (\emph{multichannel analyzer}). \\ \hline
MIO & Entradas y salidas multiplexadas del PS (\emph{Multiplexed I/O}). \\ \hline
MMCM & \emph{Mixed-Mode Clock Manager}. \\ \hline
MMU & Unidad de gestión de memoria (\emph{Memory Management Unit}). \\ \hline
MS/s, MSPS & Millones de muestras por segundo. \\ \hline
MT & Multitrigger. \\ \hline
NFS & Sistema de archivos en red (\emph{Network File System}). \\ \hline
NIM & \emph{Nuclear Instrumentation Module}. \\ \hline
OCM & Memoria en chip (\emph{On-Chip Memory}). \\ \hline
OS & Sistema operativo (\emph{operating system}). \\ \hline
OSC & Osciloscopio. \\ \hline
PC & Computadora personal. \\ \hline
PDM & Modulación por densidad de pulsos (\emph{Pulse Density Modulation}). \\ \hline
PID & Controlador proporcional-integral-derivativo. \\ \hline
PL & Lógica programable del Zynq (\emph{Programmable Logic}). \\ \hline
PLL & Lazo de enganche de fase (\emph{Phase-Locked Loop}). \\ \hline
PS & Sistema de procesamiento del Zynq (\emph{Processing System}); en algunos pasajes se lo llama \textmu{}C. \\ \hline
PSD & Discriminación por forma de pulso (\emph{Pulse Shape Discrimination}). \\ \hline
RAM & Memoria de acceso aleatorio (\emph{Random Access Memory}). \\ \hline
ROI & Región de interés (\emph{region of interest}). \\ \hline
RTL & Nivel de transferencia de registros (\emph{Register Transfer Level}). \\ \hline
SCA & Analizador monocanal (\emph{Single Channel Analyzer}). \\ \hline
SCU & Unidad de control de coherencia de caché (\emph{Snoop Control Unit}). \\ \hline
SD & Tarjeta de memoria \emph{Secure Digital}. \\ \hline
SIMD & Una instrucción, múltiples datos (\emph{Single Instruction, Multiple Data}). \\ \hline
SMA & Conector coaxial \emph{SubMiniature version A}. \\ \hline
SMP & Multiprocesamiento simétrico (\emph{Symmetric Multi-processing}). \\ \hline
SNR & Relación señal a ruido (\emph{signal-to-noise ratio}). \\ \hline
SoC & Sistema en un chip (\emph{System on Chip}). \\ \hline
SOF & Inicio de trama (\emph{start of frame}). \\ \hline
SPI & \emph{Serial Peripheral Interface}. \\ \hline
SR & Puesta a uno o a cero de un registro (\emph{set/reset}). \\ \hline
SRAM & Memoria RAM estática (\emph{Static RAM}). \\ \hline
STA & Análisis estático de tiempos (\emph{static timing analysis}). \\ \hline
SW & Software. \\ \hline
TCP & Protocolo de control de transmisión (\emph{Transmission Control Protocol}). \\ \hline
TOP & Módulo de integración: archivo cabecera del diseño. \\ \hline
TTL & Lógica transistor-transistor (\emph{Transistor-Transistor Logic}). \\ \hline
UART & Transmisor-receptor asíncrono universal (\emph{Universal Asynchronous Receiver-Transmitter}). \\ \hline
USB & \emph{Universal Serial Bus}. \\ \hline
XADC & Conversor analógico-digital interno del Zynq. \\ \hline
XDC & Archivo de restricciones de diseño de Xilinx (\emph{Xilinx Design Constraints}). \\ \hline
\textmu{}C & Microcontrolador. En el texto también designa al procesador del SoC (el PS), aunque no son necesariamente lo mismo: el PS es un procesador de aplicaciones con MMU y sistema operativo, no un microcontrolador. \\ \hline
\end{longtable}
}
